\begin{proposition}\label{B:four-by-four-five-outside}
A finite monoid with a $4\times4$ minimal band and exactly five
outside elements cannot have type $\CP^2$ or $S^2\times S^2$.
\end{proposition}
\begin{proof}
Here $A=\Q M/\Q K$ has dimension five and its two augmentation
modules $V,W$ both have dimension three.  The local case is
excluded by Proposition~\ref{B:local-three-by-three-product} for the
product and Proposition~\ref{B:local-quotient-CP2} for the projective
plane.  The residue algebra is $\Q$, since its dimension divides
both five and three.  In the nonlocal case,
Proposition~\ref{B:five-nonlocal-exceptions} leaves only the line
and Nakayama algebras.

The Nakayama Tor table in Lemma~\ref{B:five-exceptional-targets}
has no dimension-$(3,3)$ pair with Tor dimensions
$(1\text{ or }2,0,1)$.  Indeed write $v_X,w_X$ for indecomposable
multiplicities, using that lemma's notation.  Degree-two Tor gives
$(w_{T_2}+w_R)(v_{S_2}+v_L)=1$.
If $w_R=1$, degree-one vanishing forces $v_L=v_{S_1}=w_{T_1}=0$;
there is no permitted one-dimensional complement to $R$ in $W$.
Otherwise $w_{T_2}=1$, and degree-one vanishing and dimensions force
$V=S_2\oplus Ae_2$, $W=T_2\oplus e_2A$.  Their degree-zero
Tor has dimension four, again impossible.

For the line algebra write $L_a=Ae_2$, $L_b=Ae_3$ and
$R_a=e_1A$, $R_b=e_2A$.  The degree-one and degree-two conditions
leave the following possible degree-zero dimensions:
\[
\begin{array}{c|ccc}
 &T_1\oplus2T_3&T_1\oplus R_a&T_1\oplus R_b\\\hline
2S_1\oplus S_3&4&4&2\\
S_3\oplus L_a&2&1&1\\
S_3\oplus L_b&4&0&1
\end{array}
\]
These entries follow directly from the two line resolutions in
Lemma~\ref{B:five-exceptional-targets}.
In the two product entries, every balanced form has a common
one-dimensional factor and hence has rank at most one.
Lemma~\ref{B:rank-one-square} and polarization kill all products.
The middle-right projective-plane entry likewise has a rank-one
balanced form.

It remains to exclude $(V,W)=(S_3\oplus L_a,T_1\oplus R_a)$
and its opposite $(S_3\oplus L_b,T_1\oplus R_b)$.
For the first pair the invisible radical arrow is $b$.
The kernel of the action on $V$ is exactly $\Q b$; the kernel
on $W$ is the two-dimensional ideal $Ae_3A$.
Apply the complete forty-eight-basis certificate of
Lemma~\ref{B:five-line-bases}, now to $b$.
In twenty-five cases it is a basis element, so both set maps are
constant and its profile lies in the band.
In the other twenty-three it is a difference $q-p$.
The row maps are nonconstant and equal: a basis element in
$\Q b$ would make $p,q$ linearly dependent.
Fourteen cases also have nonconstant column maps, which are
therefore equal.

The other nine have $p,q\in Ae_3A$, and the following literal
left-product witnesses force their constant column maps equal:
\[
\begin{array}{c|rrrr}
\text{record}&p&q&r&t\quad(rp=rq=t)\\\hline
0&1&2&1&1\\2&1&2&1&1\\4&3&4&3&3\\5&1&2&1&1\\
18&2&3&2&2\\21&2&3&2&2\\23&2&3&2&2\\
30&2&3&1&2\\32&2&3&1&2
\end{array}
\]
Here the multiplication convention gives
$\tau_p\tau_r=\tau_t=\tau_q\tau_r$.
Thus the action-profile quotient kills $b$ modulo the band in
all cases.  Its contracted algebra is a quotient of
$A/(b)\cong\Q\times\Q(1\to2)$; every such quotient is hereditary,
since there is at most one cross arrow and no path of length two.
Its natural $H_2$ is unchanged.  The square
therefore pulls back from zero fourth homology.
The opposite argument excludes the other rank-two pair.

The program \texttt{b\_line\_4x4\_cp2.py} constructs the kernel
ideal intrinsically: choose $x$ with characters $(0,0,1)$, put
$e_3=3x^2-2x^3$, and span $Ae_3A$.  This is idempotent because
the radical squares to zero.  The independent Fraction-arithmetic
checker recomputes the ideal and every witness in both orientations,
verifying all ninety-six table/pair cases.  Hence the multiplicative
basis issue is explicitly resolved in this application as well.
\end{proof}

\begin{corollary}\label{B:fourdim-lower22}
Every finite monoid of type $\CP^2$ or $S^2\times S^2$ has at
least twenty-two elements.
\end{corollary}
\begin{proof}
The previous bounds leave only order twenty-one.
The two-row, three-by-three and three-by-four obstructions, and
the five-outside-elements bound, leave only minimal bands of
shape $3\times5$ or $4\times4$, up to opposite.
Nontrivial minimal groups require too many elements after those
shape exclusions.  Theorem~\ref{B:three-by-five-targets} handles
the first shape and Proposition~\ref{B:four-by-four-five-outside}
handles the second.
\end{proof}
